\documentclass[10pt,a4paper]{amsart}
\usepackage[margin=19mm]{geometry}
\usepackage{amsmath,amssymb,mathtools}
\usepackage[scr=rsfs,cal=euler]{mathalfa}
\usepackage{xfrac}
\usepackage[unicode,hidelinks,bookmarks=false]{hyperref}
\newcommand{\ie}{i.e.\ }
\newcommand{\cf}{cf.\ }
\newcommand{\resp}{respectively\ }
\newcommand{\Subsection}{\subsection}
\providecommand{\nobreakdash}{\nobreak\hbox{-}}
\setlength{\emergencystretch}{2em}
\multlinegap0pt
\usepackage{bm}
\usepackage{bbm}
\usepackage{dsfont}
\usepackage{xspace}
\usepackage{cancel}
\usepackage{mathtools}
\usepackage{amsthm}
\makeatletter
\@ifundefined{theorem}{\newtheorem{theorem}{Theorem}}{}
\@ifundefined{algorithm}{\newtheorem{algorithm}[theorem]{Algorithm}}{}
\makeatother
\makeatletter
\expandafter\ifx\csname proof\endcsname\relax
\newenvironment{proof}[1][Proof]{\noindent\textbf{#1.} }
{\hfill \ \rule{0.5em}{0.5em}}
\fi
\makeatother
\title[Forbidden subgraphs and complete partitions]{Forbidden subgraphs and complete partitions}
\author{John Byrne, Michael Tait,, Craig Timmons}
\date{Source: arXiv:2308.16728v2 (CC BY 4.0)}
\begin{document}
\maketitle

\noindent\emph{Source and licence.} Forbidden subgraphs and complete partitions. Version arXiv:2308.16728v2 (CC BY 4.0), distributed under the Creative Commons Attribution 4.0 International licence, as recorded in the arXiv licence metadata for this exact version and shown on the abstract page https://arxiv.org/abs/2308.16728v2. Full rights notice: licenses/arxiv-2308.16728-CC-BY-4.0.txt.\par\smallskip

\noindent\textit{Editorial scope.} Counted unit: the Algorithm whose source label is \texttt{algorithm description} in the article (source0.tex lines 918--925, source label \texttt{algorithm description}), delivered together with the article's own complete original proof of it (source0.tex lines 927--963, one \texttt{proof} environment of 37 lines). No mathematics is written, repaired or re-derived: statement and proof are the article's own text line for line. Required context retained uncounted: algorithm works (lines 289--291). Every definition, display and label of those lines is retained unchanged. Omitted from this package and not redistributed: 1--288, 292--917, 926--926, 964--1296. Nothing inside a retained excerpt is omitted or altered: every retained body is delivered exactly as the article's lines are written, so no omission receipt of any kind accompanies this package. Citations: the retained excerpts carry no citation command at all, so no bibliography entry of the article is transcribed and no reference is redistributed. Reference adaptation: none was needed and none was made; every reference inside the delivered unit resolves inside it, so no unresolved reference or citation is produced. Printed slips kept literal: no printed word, symbol, hypothesis or number of the counted statement or of its proof was altered. Layout and class: the article's own class is \texttt{article} with packages including amsmath, enumerate, amsfonts, amssymb, comment, color; the wrapper is the lane's pinned \texttt{amsart} editorial wrapper at 10pt on A4, which restates the theorem-like environments the retained text needs and adds the article's own macro definitions that the retained text needs (none), with extra wrapper packages (bm, bbm, dsfont, xspace, cancel, mathtools). The delivered numbering is the unit's own. Assets: no retained span contains a figure environment, a graphics-inclusion command, a PDF-page inclusion, a table or a TikZ picture; no image file is copied into this package, the package declares no asset dependency and no image file is redistributed with it. Licence: version arXiv:2308.16728v2 of this article is distributed under the Creative Commons Attribution 4.0 International licence, as recorded in the arXiv licence metadata for this exact version and shown on the abstract page https://arxiv.org/abs/2308.16728v2. A full rights notice is delivered at licenses/arxiv-2308.16728-CC-BY-4.0.txt. Characters: every retained excerpt is pure ASCII, so the delivered engine, XeLaTeX with its default Latin Modern text font, reports no missing character.
\begin{theorem} \label{algorithm works}
    Let $H$ be a bipartite graph with no vertices of degree 1. Suppose there exists an $H$-free $d$-regular graph $G$ on $n$ vertices, where $d=n^a$ and $a<\frac{1}{3}$. Let $\rho = \max\{\rho_2, -\rho_n\}$ if $G$ is not bipartite and $\rho = \rho_2$ if $G$ is bipartite and assume that $\rho=O(d^\frac{1}{2})$. Then Algorithm \ref{algorithm description} terminates with a complete partition of a graph $G'$, with $m=\frac{1}{4}n^{\frac{1+a}{2}}/(\log n)^{\frac{1}{2}}$ parts and maximum part size at most $Cm^{\frac{2}{1+a}-1}(\log m)^{\frac{1}{1+a}}$, where $C$ is an absolute constant.
\end{theorem}

\begin{algorithm} \label{algorithm description} Let $H$ and $G$ be graphs satisfying the assumptions of Theorem \ref{algorithm works}.
    \begin{enumerate}
        \item Initialize $V_1=\cdots=V_m=\emptyset$.
        \item For each $i$, add $n^{\frac{1-a}{2}}(\log n)^{\frac{1}{2}}$ vertices in $G$ to $V_i$ such that each vertex added is not already contained in any other $V_j$ and is also at distance at least $3$ from every other vertex in $V_i$. If $G$ is bipartite then all added vertices should also belong to $X$.
        \item For $1\le i\le m$, define $\mathcal S_i=\{j:j\ne i,e(V_i,V_j)=0\}$ and $s_i=|\mathcal S_i|$. If for every $i$ we have $s_i<n^{\frac{1-a}{2}}(\log n)^{\frac{1}{2}}$, proceed to step 4. Otherwise, for each $1\le i\le m$, choose the vertex in $B-(N(V_i)\cup N(N(V_i)))-(V_1\cup\cdots\cup V_m)$ which is adjacent to the maximum number of vertices in $\bigcup_{j\in\mathcal S_i}V_j$ and add it to $V_i$. Return to the beginning of step 3.
        \item For each $i$, for each $j\in\mathcal S_i$, create a new vertex in $V_i$ and make it adjacent to some vertex in $V_j$ which was added to $V_j$ in one of the previous steps.
\end{enumerate}
\end{algorithm}

\begin{proof}[Proof of Theorem \ref{algorithm works}] Is it possible to perform step 2, because at any point during step 2 when we want to add a vertex to $V_i$, we have
$$
\begin{aligned}
    |A-(N(V_i)\cup N(N(V_i)))-(V_1\cup\cdots\cup V_m)|&\ge\frac{n}{2}-d^2|V_i|-m\cdot|V_1|\\ &\ge \frac{n}{2}-n^{2a+\frac{1-a}{2}}(\log n)^{\frac{1}{2}}-\frac{1}{4}n^{\frac{1+a}{2}+\frac{1-a}{2}}.
\end{aligned}$$
The choice of $a$ ensures that both exponents in the right-hand side are at most 1 and the right-hand side is positive for $n$ sufficiently large.

To show it is possible to perform step 3, we track how the number $s_i$ changes each time a vertex is added. 
We assume that $i$ is a value for which the inequality 
$s_i \geq n^{ \frac{1-a}{2} } ( \log n )^{1/2}$
is true.  
Let $s_i(t)$ be the value of $s_i$ after $t$ iterations of step 3 have been performed. Now during any of the first $2an^{\frac{1-a}{2}}(\log n)^{\frac{1}{2}}$ iterations of step 3 we have
$$\begin{aligned}|B-(N(V_i)\cup N(N(V_i)))-(V_1\cup\cdots\cup& V_m)|\ge  \frac{n}{2}-n^{2a}(n^\frac{1-a}{2}(\log n)^{\frac{1}{2}}+2an^{\frac{1-a}{2}}(\log n)^{\frac{1}{2}})\\
&-\frac{1}{4}\frac{n^{\frac{1+a}{2}}}{(\log n)^{\frac{1}{2}}}(n^\frac{1-a}{2}(\log n)^{\frac{1}{2}}+2an^{\frac{1-a}{2}}(\log n)^{\frac{1}{2}}).\end{aligned}$$
The choice of $a$ ensures that the right-hand side is $\Theta(n)$. Set $W=B-(N(V_i)\cup N(N(V_i)))-(V_1\cup\cdots\cup V_m)$. To apply the expander mixing lemma we would like
$$\frac{d\left|\bigcup_{j\in\mathcal S_i(t)}V_j\right||W|}{n}\gg\rho\sqrt{\left|\bigcup_{j\in\mathcal S_i(t)}V_j\right||W|},$$
in other words $\Theta(n)=|W|\gg\frac{\rho^2n^2}{d^2\left|\bigcup_{j\in\mathcal S_i(t)}V_j\right|}$. Noting that 
$\rho= O (d^{\frac{1}{2}})$ 
and $\left|\bigcup_{j\in\mathcal S_i(t)}V_j\right|\ge s_i(t)n^{\frac{1-a}{2}}(\log n)^{\frac{1}{2}}\ge n^{1-a}(\log n)$ we see that this condition is met. Therefore,
$$e\left(W,\bigcup_{j\in\mathcal S_i(t)}V_j\right)\ge\frac{1}{2}\frac{ds_i(t)n^{\frac{1-a}{2}}(\log n)^{\frac{1}{2}}|W|}{n}=\frac{1}{2n^\frac{1-a}{2}}|W|s_i(t)(\log n)^{\frac{1}{2}}.$$
Thus, there exists a vertex in $W$ with at least $\frac{1}{2n^\frac{1-a}{2}}s_i(t)(\log n)^{\frac{1}{2}}$ neighbors in $\left|\bigcup_{j\in\mathcal S_i(t)}V_j\right|$, and so in iteration $t$ of step 3, a vertex is added to $V_i$ which has at least $\frac{1}{2n^{\frac{1-a}{2}}}s_i(t)(\log n)^{\frac{1}{2}}$ neighbors in $\left|\bigcup_{j\in\mathcal S_i(t)}V_j\right|$. Now since each of the sets $V_j$ has all of its vertices at distance at least 3 from each other, each neighbor of the vertex added belongs to a different set $V_j$ (where $j\in\mathcal S_i(t)$). It follows that
$$s_i(t+1)\le s_i(t)-\frac{1}{2n^\frac{1-a}{2}}(\log n)^{\frac{1}{2}}s_i(t).$$
Using $s_i(0)\le m\le n^\frac{1+a}{2}/(\log n)^{\frac{1}{2}}$, we then have that as long as $|W|$ is large enough as above, 
$$s_i(t)\le \frac{n^{\frac{1+a}{2}}}{(\log n)^{\frac{1}{2}}}\left(1-\frac{1}{2n^\frac{1-a}{2}}(\log n)^{\frac{1}{2}}\right)^t.$$
Using the estimate $\log(1-x)\le -x$, we therefore have
$$\log s_i(t)\le\left(\frac{1+a}{2}\log n-\frac{1}{2}\log\log n\right)-t\frac{(\log n)^{\frac{1}{2}}}{2n^{\frac{1-a}{2}}}.$$
Step 3 terminates when for all $i$, $s_i(t) < n^{\frac{1-a}{2}}(\log n)^{\frac{1}{2}}$. Solving for $t$ in the inequality
$$\left(\frac{1+a}{2}\log n-\frac{1}{2}\log\log n\right)-t\frac{(\log n)^{\frac{1}{2}}}{2n^{\frac{1-a}{2}}}\le\log\left(n^{\frac{1-a}{2}}(\log n)^{\frac{1}{2}}\right)$$
we see that this is guaranteed to happen by iteration $t=2a(\log n)^{\frac{1}{2}}n^{\frac{1-a}{2}}.$

In step 4, none of the vertices created creates a copy of $H$, since $H$ has no vertex of degree 1. 

We calculate the maximum size of a part in the resulting $(m,k)$-graph. For each $V_i$, $n^{\frac{1-a}{2}}(\log n)^{\frac{1}{2}}$ vertices were added in step 2, at most $2an^{\frac{1-a}{2}}(\log n)^{\frac{1}{2}}$ were added in step 3, and at most $n^{\frac{1-a}{2}}(\log n)^{\frac{1}{2}}$ were added in step 4. Therefore we have $f(m,H)\le 4n^{\frac{1-a}{2}}(\log n)^{\frac{1}{2}}$. Choose $C$ such that $C\left(\frac{1}{4}\right)^{\frac{2}{1+a}-1}\left(\frac{1}{2}\right)^{\frac{1}{1+a}}\ge 4$. Using the fact that
$$\log m=\frac{1+a}{2}\log n-\log4-\frac{1}{2}\log\log n\ge\frac{1}{2}\log n$$
for $n$ large enough, we find that
$$Cm^{\frac{2}{1+a}-1}(\log m)^{\frac{1}{1+a}}\ge 4n^{\frac{1-a}{2}}(\log n)^{\frac{1}{2}}\ge f(m,H).$$
\end{proof}

\end{document}
